\section{Literature Review}
\label{sec:literature}

\subsection{DARP, DRT, and On-Demand Mobility Operations}

The Dial-a-Ride Problem (DARP) models passenger transport with origin--destination
requests, time windows, ride-time limits, vehicle capacities, and fleet routing decisions
\citep{cordeau2007,parragh2008a}. It has long provided an operations-research foundation
for paratransit and shared passenger transport, and more recent work connects DARP
methods to demand-responsive transit (DRT), ridepooling, and on-demand mobility
logistics \citep{ho2018,molenbruch2017}. In these systems, the operator must manage a
fleet under sparse and dynamic request patterns, making routing intensity, service
coverage, and operational reliability central logistics and operations concerns.

Classical DARP assumes door-to-door service. This maximizes passenger convenience but
can fragment vehicle tours when many-to-many requests are spatially dispersed. Meeting
points or virtual stops introduce service consolidation into the passenger-transport
setting: passengers walk a bounded distance so that vehicles can serve more compatible
stops and reduce detours. This paper treats meeting-point DRT as an operational
service-design problem, not only as a routing variant, because the consolidation choice
changes both fleet operations and passenger willingness to accept the offered service.

\subsection{Meeting-Point DRT and Service Consolidation}

Meeting points were introduced into shared-ride and ridepooling research as a way to
trade small walking burdens for route consolidation \citep{stiglic2015,stiglic2018}.
In the DARP literature, \citet{cortenbach2024} formalize the DARP with meeting points
(DARPmp), allowing pickup-side flexibility while retaining fixed dropoff destinations.
That work demonstrates the operational value of flexible pickup locations, but its
single-sided design leaves dropoff-side consolidation outside the model and assumes that
assigned passengers comply with the service offer.

Bidirectional walking flexibility has also been studied in ridepooling contexts.
\citet{fielbaum2021} show that allowing both pickup and dropoff walking can reduce
vehicle detours in static shared-ride settings. The present paper transfers that
service-consolidation idea to many-to-many DRT operations with online request arrivals,
rolling-horizon routing, and simulated passenger choice. The bidirectional design is
therefore positioned as a logistics and operations service-design lever: it may lower
routing intensity for served trips, but it can also alter served share and rejection
composition.

\subsection{Passenger Choice in Service Design}

Discrete choice models provide a standard framework for representing heterogeneous
passenger preferences over service attributes such as walking, waiting, in-vehicle time,
and price \citep{mcfadden1974,benakiva1985,train2009,lavieri2019}. Most DRT routing
models represent rejection as a feasibility outcome rather than an acceptance response:
if a request can be inserted within the constraints, it is treated as served. That
deterministic treatment is limiting for meeting-point service designs because walking
burden and service delay are precisely the attributes that may cause passengers to reject
an offer.

This paper uses a single-offer binary-logit mechanism as a simulation construct. The
operator presents one routing-feasible bundle, and the simulated passenger either accepts
that bundle or takes an outside option. Passenger types are used to represent a range of
service sensitivities within the simulation; they are not empirical population segments
validated from field data. This choice layer makes it possible to distinguish
choice rejection from feasibility rejection and to evaluate passenger-segment monitoring
implications without claiming real population equity validation.

\subsection{Dynamic DRT Routing and Rolling-Horizon Operations}

DRT is commonly operated online: requests arrive over time, vehicles move while new
offers are being made, and dispatch plans must be updated repeatedly. Dynamic vehicle
routing and online DARP research therefore use insertion heuristics, metaheuristics, and
rolling-horizon approaches to keep decisions tractable under uncertainty
\citep{psaraftis1988,pillac2013,ho2018}. Rolling-horizon methods solve a finite active
window, commit near-term decisions, and then reoptimize as new information arrives
\citep{bent2004}. This structure is well suited to DRT operations because it represents
the practical rhythm of dispatch rather than a clairvoyant offline plan.

Adaptive Large Neighborhood Search (ALNS) is widely used for DARP and related routing
problems because destroy-and-repair operators can search large neighborhoods while
respecting feasibility constraints \citep{ropke2006}. \citet{wu2025} apply rolling-horizon
ALNS to many-to-many DRT, showing the operational importance of periodic
reoptimization. Their setting remains door-to-door and deterministic with respect to
passenger acceptance. The present paper combines rolling-horizon operations with
bidirectional meeting-point service consolidation and binary-logit passenger response.

\subsection{Research Gap and Positioning}

Table~\ref{tab:literature-gap} summarizes the positioning of this paper relative to the
most closely related prior works. The contribution is a passenger-response-aware
simulation framework for bidirectional meeting-point DRT service design: bidirectional
pickup and dropoff meeting points, single-offer binary-logit response, and
rolling-horizon many-to-many operations are evaluated together while keeping diagnostic
evidence roles separate from the primary behavioral narrative.

\begin{table}[h]
\centering
\caption{Positioning relative to prior work}
\label{tab:literature-gap}
\begin{tabular}{lcccc}
\toprule
Study & Bidirectional MPs & Passenger Choice & Rolling Horizon & Many-to-Many \\
\midrule
Cortenbach et al.\ (2024) & No (single-sided) & No & No & Yes \\
Wu et al.\ (2025)         & No                & No & Yes & Yes \\
Fielbaum et al.\ (2021)   & Yes (ridepooling) & No & No & Partial \\
Alonso-Mora et al.\ (2017) & No               & No & No & Yes \\
This paper                & Yes               & Yes (binary logit) & Yes & Yes \\
\bottomrule
\end{tabular}
\end{table}
